% preamble.tex - Documento definitivo (unificado)
% Base: preámbulo del proyecto Overleaf (bibtex IEEEtran + cleveref + algorithm2e)
% Añadidos: TikZ, multirow y placeins, requeridos por la Introducción y el Marco Teórico.

\usepackage[spanish,es-minimal]{babel}
\usepackage[utf8]{inputenc}
\usepackage[T1]{fontenc}
\usepackage{lmodern}
\usepackage{microtype}

% Matemáticas
\usepackage{amsmath, amssymb, amsthm}

% Gráficos y tablas
\usepackage{graphicx}
\usepackage{booktabs}
\usepackage{longtable}
\usepackage{array}
\usepackage{tabularx}
\usepackage{multirow}
\usepackage{float}
\usepackage{placeins}

% Diagramas (Introducción y Marco Teórico)
\usepackage{tikz}
\usetikzlibrary{arrows.meta, positioning, shapes.geometric, calc, fit}

% Maquetación
\usepackage[margin=2.5cm]{geometry}
\usepackage{setspace}
\usepackage{parskip}
\usepackage{enumitem}

% Referencias cruzadas
\usepackage{cite}        % agrupa/ordena citas numéricas estilo IEEE: [1]-[3]
\usepackage{hyperref}
\usepackage{cleveref}

% --- Bibliografía IEEE con BibTeX clásico (rápido; reemplaza a biblatex/biber) ---
% El estilo IEEEtran.bst numera por orden de aparición (equivalente a sorting=none).
\renewcommand{\bibname}{Referencias}

\usepackage[font=small,labelfont=bf]{caption}
\captionsetup[figure]{name=Fig.,labelsep=period}
\captionsetup[table]{name=TABLA,labelsep=newline,position=above,
  font=small,labelfont={bf,sc},textfont=sc}

\usepackage{chngcntr}
\counterwithin{figure}{chapter}
\counterwithin{table}{chapter}
\counterwithin{equation}{chapter}

% Pseudocódigo (estilo IEEE)
\usepackage[ruled,vlined,linesnumbered,spanish]{algorithm2e}
\SetAlgorithmName{Algoritmo}{algoritmo}{Lista de algoritmos}

% Etiquetas de cleveref en español
\crefname{figure}{Figura}{Figuras}
\crefname{table}{Tabla}{Tablas}
\crefname{equation}{Ecuación}{Ecuaciones}
\crefname{section}{Sección}{Secciones}
\crefname{chapter}{Capítulo}{Capítulos}
\crefname{algocf}{Algoritmo}{Algoritmos}
\Crefname{algocf}{Algoritmo}{Algoritmos}

% Color para énfasis puntual
\usepackage{xcolor}

% Configuración de hyperref
\hypersetup{
    colorlinks=true,
    linkcolor=black!60!blue,
    citecolor=black!60!blue,
    urlcolor=black!60!blue,
}

% Rutas de figuras
\graphicspath{{./}{./figures/}{./figures/ch09/}}
